\label{kap:statistika}
\mathquote{There are three kinds of lies: lies, damned lies, and statistics.}{Mark Twain (který to připsal Benjaminu Disraelimu, který to nejspíš nikdy neřekl)}

Pravděpodobnost se ptá: když víme, jak funguje náhoda, co uvidíme? Statistika se ptá obráceně: když něco vidíme, co to říká o světě? V této kapitole se podíváme na \textbf{popisnou statistiku}, tedy na to, jak velké množství dat shrnout několika čísly a grafy. Je to ta část matematiky, kterou nejčastěji uvidíte ve zprávách, a také ta, se kterou se nejčastěji manipuluje.

\section{Základní pojmy}\label{sec:stat-zakladni-pojmy}

\begin{definitionbox}
  \begin{itemize}
    \item \textbf{Statistický soubor} je množina objektů, které zkoumáme (žáci třídy, obyvatelé města, výrobky z linky).
    \item \textbf{Statistická jednotka} je jeden prvek souboru. Počet jednotek je \textbf{rozsah souboru} $n$.
    \item \textbf{Statistický znak} je vlastnost, kterou u jednotek sledujeme. Může být \textbf{kvantitativní} (číselný: výška, známka, plat) nebo \textbf{kvalitativní} (slovní: barva očí, oblíbený předmět).
    \item \textbf{Absolutní četnost} $n_i$ hodnoty $x_i$ je počet jednotek, které mají tuto hodnotu. \textbf{Relativní četnost} je $\frac{n_i}{n}$ (často v procentech). \textbf{Kumulativní četnost} je součet četností všech hodnot menších nebo rovných $x_i$.
  \end{itemize}
\end{definitionbox}

Často nezkoumáme celý soubor (\textbf{základní soubor}, například všechny voliče), ale jen jeho část (\textbf{výběrový soubor}, například tisíc dotázaných v průzkumu). Aby výsledky něco znamenaly, musí být výběr \textbf{reprezentativní}. Slavný případ opaku: v roce 1936 americký časopis \textit{Literary Digest} obeslal přes deset milionů lidí a předpověděl drtivé vítězství republikána Landona nad Rooseveltem. Roosevelt vyhrál ve 46 státech ze 48. Časopis totiž vybíral adresy z telefonních seznamů a seznamů majitelů aut, a ty si uprostřed hospodářské krize mohli dovolit hlavně bohatší voliči. O dva roky později časopis zanikl.

\begin{example}\label{ex:znamky}
  Známky z písemky v malé skupině $10$ studentů: $1, 2, 2, 3, 1, 4, 2, 3, 5, 2$. Tabulka četností:
  \begin{center}
    \renewcommand{\arraystretch}{1.3}
    \begin{tabular}{|c|c|c|c|c|c|}
      \hline
      známka $x_i$ & $1$ & $2$ & $3$ & $4$ & $5$ \\
      \hline
      absolutní četnost $n_i$ & $2$ & $4$ & $2$ & $1$ & $1$ \\
      relativní četnost & $20\,\%$ & $40\,\%$ & $20\,\%$ & $10\,\%$ & $10\,\%$ \\
      kumulativní četnost & $2$ & $6$ & $8$ & $9$ & $10$ \\
      \hline
    \end{tabular}
  \end{center}
\end{example}

\section{Grafické znázornění}\label{sec:graficke-znazorneni}

\begin{definitionbox}
  \begin{itemize}
    \item \textbf{Sloupcový diagram} (histogram): výška sloupce odpovídá četnosti. Vhodný pro kvantitativní znaky.
    \item \textbf{Kruhový (koláčový) diagram}: velikost výseče odpovídá relativní četnosti. Úhel výseče je $\frac{n_i}{n}\cdot 360^\circ$. Vhodný pro podíly celku.
    \item \textbf{Spojnicový graf}: vhodný pro vývoj v čase.
    \item \textbf{Krabicový graf} (box plot): zobrazuje minimum, kvartily a maximum, viz dále.
  \end{itemize}
\end{definitionbox}

\begin{center}
  \begin{tikzpicture}[scale=0.8]
    \draw[->] (0,0) -- (6.2,0) node[right] {\small známka};
    \draw[->] (0,0) -- (0,4.8) node[above] {\small četnost};
    \foreach \y in {1,2,3,4} \draw (0.08,\y) -- (-0.08,\y) node[left] {\tiny $\y$};
    \foreach \x/\h in {1/2, 2/4, 3/2, 4/1, 5/1} {
      \fill[colorofwisdom!60] (\x-0.35,0) rectangle (\x+0.35,\h);
      \draw (\x-0.35,0) rectangle (\x+0.35,\h);
      \node[below] at (\x,0) {\small $\x$};
    }
  \end{tikzpicture}
  \hspace{1.5cm}
  \begin{tikzpicture}[scale=1.3]
    \fill[colorofwisdom!80] (0,0) -- (0:1) arc (0:72:1) -- cycle;
    \fill[colorofwisdom!40] (0,0) -- (72:1) arc (72:216:1) -- cycle;
    \fill[black!30] (0,0) -- (216:1) arc (216:288:1) -- cycle;
    \fill[black!60] (0,0) -- (288:1) arc (288:324:1) -- cycle;
    \fill[black!85] (0,0) -- (324:1) arc (324:360:1) -- cycle;
    \draw (0,0) circle (1);
    \node at (36:0.65) {\small $1$};
    \node at (144:0.65) {\small $2$};
    \node at (252:0.65) {\small $3$};
    \node[white] at (306:0.7) {\small $4$};
    \node[white] at (342:0.7) {\small $5$};
  \end{tikzpicture}
\end{center}

\begin{note}
  Grafy jsou nejsnadnější způsob, jak se statistikou lhát, aniž byste řekli jedinou nepravdu. Oblíbené triky: osa $y$, která nezačíná nulou (malý rozdíl pak vypadá obrovsky), trojrozměrné koláčové grafy (přední výseče vypadají větší) a vybírání časového úseku tak, aby trend vypadal, jak se nám hodí. Když vidíte graf ve zprávách, nejdřív se podívejte na osy.

  \begin{center}
    \begin{tikzpicture}[scale=0.8]
      \begin{scope}
        \draw[->] (0,0) -- (3.4,0);
        \draw[->] (0,0) -- (0,3.4);
        \foreach \y/\l in {1/20, 2/40} \draw (0.07,\y) -- (-0.07,\y) node[left] {\tiny $\l\,\%$};
        \filldraw[fill=gray!40, draw=black] (0.5,0) rectangle (1.4,2.5);
        \filldraw[fill=colorofwisdom!70, draw=black] (2,0) rectangle (2.9,2.6);
        \node[below, font=\scriptsize] at (0.95,0) {strana A};
        \node[below, font=\scriptsize] at (2.45,0) {strana B};
        \node[above, font=\scriptsize] at (0.95,2.5) {$50\,\%$};
        \node[above, font=\scriptsize] at (2.45,2.6) {$52\,\%$};
        \node[font=\small] at (1.7,-1) {osa od nuly};
      \end{scope}
      \begin{scope}[xshift=6cm]
        \draw[->] (0,0) -- (3.4,0);
        \draw[->] (0,0) -- (0,3.4);
        \draw (-0.12,0.15) -- (0.12,0.25);
        \draw (-0.12,0.25) -- (0.12,0.35);
        \foreach \y/\l in {0.9/50, 2.6/52} \draw (0.07,\y) -- (-0.07,\y) node[left] {\tiny $\l\,\%$};
        \filldraw[fill=gray!40, draw=black] (0.5,0) rectangle (1.4,0.9);
        \filldraw[fill=colorofwisdom!70, draw=black] (2,0) rectangle (2.9,2.6);
        \node[below, font=\scriptsize] at (0.95,0) {strana A};
        \node[below, font=\scriptsize] at (2.45,0) {strana B};
        \node[font=\small] at (1.7,-1) {osa od $49\,\%$};
      \end{scope}
    \end{tikzpicture}
  \end{center}
  Oba grafy ukazují stejná data. V levém je vidět, že strany jsou skoro vyrovnané. V pravém to vypadá, že strana B vede trojnásobně, protože osa začíná těsně pod $50\,\%$.
\end{note}

\section{Charakteristiky polohy}\label{sec:charakteristiky-polohy}

Charakteristiky polohy odpovídají na otázku `jaká je typická hodnota?'.

\begin{definitionbox}
  Pro hodnoty $x_1, x_2, \ldots, x_n$:
  \begin{itemize}
    \item \textbf{Aritmetický průměr}:
    \begin{equation*}
      \bar{x} = \frac{x_1 + x_2 + \cdots + x_n}{n} = \frac{1}{n}\sum_{i=1}^{n}x_i.
    \end{equation*}
    Pokud se hodnoty $x_i$ opakují s četnostmi $n_i$, je to \textbf{vážený průměr} $\bar{x} = \frac{\sum n_ix_i}{\sum n_i}$.
    \item \textbf{Medián} $\tilde{x}$: prostřední hodnota po seřazení dat podle velikosti. Pro sudé $n$ je to průměr dvou prostředních hodnot.
    \item \textbf{Modus} $\hat{x}$: nejčastější hodnota.
    \item \textbf{Geometrický průměr} (pro kladné hodnoty): $\bar{x}_G = \sqrt[n]{x_1x_2\cdots x_n}$.
    \item \textbf{Harmonický průměr}: $\bar{x}_H = \dfrac{n}{\frac{1}{x_1} + \cdots + \frac{1}{x_n}}$.
  \end{itemize}
\end{definitionbox}

\begin{example}
  V datech z příkladu~\ref{ex:znamky} je průměr $\bar{x} = \frac{1\cdot 2 + 2\cdot 4 + 3\cdot 2 + 4 + 5}{10} = \frac{25}{10} = 2{,}5$. Seřazená data: $1, 1, 2, 2, 2, 2, 3, 3, 4, 5$. Medián je průměr pátého a šestého prvku, tedy $\tilde{x} = 2$. Modus je $\hat{x} = 2$.

  Průměr má krásný fyzikální význam: je to \textbf{těžiště}. Představte si číselnou osu jako lehké prkno a na každou hodnotu položte jednu stejně těžkou kuličku. Prkno bude v rovnováze, právě když ho podepřete v průměru. Odchylky nalevo ($-1{,}5$ dvakrát a $-0{,}5$ čtyřikrát, celkem $-5$) se přesně vyváží s odchylkami napravo ($0{,}5$ dvakrát, $1{,}5$ a $2{,}5$, celkem $+5$).
  \begin{center}
    \begin{tikzpicture}[scale=1]
      \draw[thick] (0.3,0) -- (8.7,0);
      \foreach \x in {1,...,5} \draw (1.6*\x,0.08) -- (1.6*\x,-0.08) node[below=1pt] {\small $\x$};
      \foreach \x/\n in {1/2, 2/4, 3/2, 4/1, 5/1} {
        \foreach \k in {1,...,\n} {
          \filldraw[fill=colorofwisdom!60, draw=black] (1.6*\x,{0.28*\k-0.1}) circle (0.12);
        }
      }
      \fill[colorofwisdom] (4,-0.42) -- (3.8,-0.8) -- (4.2,-0.8) -- cycle;
      \node[right, font=\small, colorofwisdom] at (4.25,-0.65) {$\bar{x} = 2{,}5$};
      \draw[dashed, gray] (3.2,-0.3) -- (3.2,1.4) node[above, black, font=\small] {$\tilde{x} = 2$};
    \end{tikzpicture}
  \end{center}
  Medián naopak dělí kuličky na dvě stejně \textit{početné} poloviny, ať jsou krajní kuličky jakkoli daleko.
\end{example}

\begin{example}
  Proč potřebujeme medián? Firma má $9$ zaměstnanců s platem $30\,000$ Kč a ředitele s platem $300\,000$ Kč. Průměrný plat je
  \begin{equation*}
    \bar{x} = \frac{9\cdot 30000 + 300000}{10} = 57\,000\ \text{Kč}.
  \end{equation*}
  Ale nikdo kromě ředitele tolik nebere! Medián je $30\,000$ Kč a mnohem lépe popisuje `typického' zaměstnance. Průměr je citlivý na extrémní hodnoty, medián ne. Proto se ve zprávách o mzdách vždy vyplatí zjistit, zda jde o průměr, nebo o medián. Medián mezd v Česku bývá zhruba o šestinu nižší než průměr.

  \begin{center}
    \begin{tikzpicture}[scale=1]
      \draw[->] (0,0) -- (11,0) node[right] {\small tis. Kč};
      \foreach \x in {0,50,...,300} \draw ({0.035*\x},0.08) -- ({0.035*\x},-0.08) node[below] {\scriptsize $\x$};
      \foreach \k in {1,...,9} \filldraw[fill=gray!50, draw=black] (1.05,{0.24*\k-0.08}) circle (0.1);
      \filldraw[fill=colorofwisdom!70, draw=black] (10.5,0.16) circle (0.1);
      \node[above, font=\scriptsize] at (10.5,0.3) {ředitel};
      \draw[dashed, gray] (1.05,-0.5) -- (1.05,2.5) node[above, black, font=\small] {medián $30$};
      \fill[colorofwisdom] (1.995,-0.45) -- (1.82,-0.8) -- (2.17,-0.8) -- cycle;
      \node[right, font=\small, colorofwisdom] at (2.2,-0.65) {průměr $57$};
    \end{tikzpicture}
  \end{center}
  V řeči těžiště: jedna kulička daleko vpravo má dlouhé `rameno páky', takže přetáhne podpěru (průměr) doprava, přestože nalevo je kuliček devět. Medián se o vzdálenost nestará, jen o počet.
\end{example}

\begin{example}\label{ex:jine-prumery}
  Kdy použít jiné průměry?
  \begin{itemize}
    \item \textbf{Harmonický průměr} pro průměrnou rychlost. Pojedete do Brna rychlostí $60\ \si{km/h}$ a zpátky rychlostí $40\ \si{km/h}$. Průměrná rychlost není $50$, ale $\frac{2}{\frac{1}{60} + \frac{1}{40}} = 48\ \si{km/h}$, protože pomalejší cestou strávíte víc času. Pokud pochybujete, spočítejte to pro konkrétní vzdálenost, třeba $120\ \si{km}$: tam $2$ hodiny, zpátky $3$ hodiny, celkem $240\ \si{km}$ za $5$ hodin.
    \item \textbf{Geometrický průměr} pro průměrný růst. Akcie vzrostla první rok o $10\,\%$ a druhý o $60\,\%$. Průměrný roční růst není $35\,\%$, ale $\sqrt{1{,}1\cdot 1{,}6} - 1 = \sqrt{1{,}76} - 1\approx 32{,}7\,\%$. Kontrola: $1{,}327^2\approx 1{,}76$.
  \end{itemize}
\end{example}

\begin{note}
  Pro kladná čísla vždy platí $\bar{x}_H\leq\bar{x}_G\leq\bar{x}$, a rovnost nastává jen tehdy, když jsou všechna čísla stejná. Pro dvě čísla $a, b$ se nerovnost $\sqrt{ab}\leq\frac{a + b}{2}$ dokáže jednoduše: je ekvivalentní s $0\leq{(\sqrt{a} - \sqrt{b})}^2$. Geometricky to znamená, že ze všech obdélníků s daným obvodem má největší obsah čtverec, což jsme viděli v kapitole~\ref{kap:elementarni}.
\end{note}

\section{Charakteristiky variability}\label{sec:charakteristiky-variability}

Dva soubory mohou mít stejný průměr a přitom vypadat úplně jinak. Data $\{4, 5, 6\}$ a $\{0, 5, 10\}$ mají oba průměr $5$, ale ty druhé jsou mnohem víc `rozházené'. Charakteristiky variability měří, jak moc se data od průměru odchylují.

\begin{center}
  \begin{tikzpicture}[scale=0.85]
    \foreach \row/\vals/\s/\lab in {0/{4,5,6}/0.816/{$\{4, 5, 6\}$: $s\approx 0{,}82$}, -1.6/{0,5,10}/4.082/{$\{0, 5, 10\}$: $s\approx 4{,}08$}} {
      \fill[colorofwisdom!12] ({5-\s},\row-0.35) rectangle ({5+\s},\row+0.35);
      \draw (-0.3,\row) -- (10.3,\row);
      \foreach \x in {0,...,10} \draw (\x,\row+0.06) -- (\x,\row-0.06);
      \foreach \v in \vals {
        \ifnum\v=5\relax\else
          \draw[->, colorofwisdom, thick] (5,\row+0.2) -- (\v,\row+0.2);
        \fi
        \filldraw[fill=colorofwisdom!60, draw=black] (\v,\row) circle (0.12);
      }
      \node[anchor=west, font=\small] at (10.6,\row) {\lab};
    }
    \foreach \x in {0,5,10} \node[below, font=\scriptsize] at (\x,-1.7) {$\x$};
    \fill[colorofwisdom] (5,-2.15) -- (4.85,-2.45) -- (5.15,-2.45) -- cycle;
    \node[right, font=\scriptsize] at (5.15,-2.3) {$\bar{x} = 5$};
  \end{tikzpicture}
\end{center}
Šipky jsou odchylky od průměru, světlý pruh sahá od $\bar{x} - s$ do $\bar{x} + s$. Směrodatnou odchylku si můžete představit jako `typickou vzdálenost' hodnoty od průměru. U horních dat je to necelá jednička, u dolních přes čtyři.

\begin{definitionbox}
  \begin{itemize}
    \item \textbf{Variační rozpětí}: $R = x_{\max} - x_{\min}$.
    \item \textbf{Rozptyl} (průměrná druhá mocnina odchylky od průměru):
    \begin{equation*}
      s^2 = \frac{1}{n}\sum_{i=1}^{n}{(x_i - \bar{x})}^2 = \overline{x^2} - \bar{x}^2.
    \end{equation*}
    \item \textbf{Směrodatná odchylka}: $s = \sqrt{s^2}$. Má stejné jednotky jako data.
    \item \textbf{Variační koeficient}: $v = \frac{s}{\bar{x}}$ (relativní míra variability, obvykle v procentech).
  \end{itemize}
\end{definitionbox}

Proč druhé mocniny? Kdybychom průměrovali přímo odchylky $x_i - \bar{x}$, vyšla by vždy nula, protože kladné a záporné odchylky se přesně vyruší (zkuste si to dokázat). Druhá mocnina udělá ze všech odchylek kladná čísla. Směrodatná odchylka pak vrací výsledek do původních jednotek. Vzorec $s^2 = \overline{x^2} - \bar{x}^2$ (průměr druhých mocnin minus druhá mocnina průměru) se počítá rychleji a odvodí se roznásobením.
 Druhá mocnina má i geometrický smysl: nad každou odchylkou si nakreslete čtverec. Rozptyl je \textbf{průměrný obsah těchto čtverců} a směrodatná odchylka je strana čtverce s tímto průměrným obsahem. Vzdálené hodnoty mají obrovské čtverce, takže rozptyl na ně reaguje mnohem citlivěji než na hodnoty blízko průměru.

\begin{example}
  Pro známky z příkladu~\ref{ex:znamky} s průměrem $2{,}5$:
  \begin{align*}
    s^2 &= \frac{2\cdot{(1 - 2{,}5)}^2 + 4\cdot{(2 - 2{,}5)}^2 + 2\cdot{(3 - 2{,}5)}^2 + {(4 - 2{,}5)}^2 + {(5 - 2{,}5)}^2}{10}\\
    &= \frac{4{,}5 + 1 + 0{,}5 + 2{,}25 + 6{,}25}{10} = 1{,}45.
  \end{align*}
  Směrodatná odchylka $s = \sqrt{1{,}45}\approx 1{,}20$ a variační koeficient $v = \frac{1{,}20}{2{,}5}\approx 48\,\%$. Kontrola druhým vzorcem: $\overline{x^2} = \frac{2 + 16 + 18 + 16 + 25}{10} = 7{,}7$ a $7{,}7 - 6{,}25 = 1{,}45$.
\end{example}

\begin{note}
  Ve statistice výběrových souborů se rozptyl často počítá s $n - 1$ místo $n$ ve jmenovateli (tzv. výběrový rozptyl), protože to dává lepší odhad rozptylu celé populace. U maturity se obvykle používá verze s $n$, ale pokud vám kalkulačka nebo tabulkový procesor vyhodí jiné číslo, je to pravděpodobně tohle.
\end{note}

\begin{definitionbox}
  \textbf{Kvantily}: \textbf{Dolní kvartil} $Q_1$ je hodnota, pod kterou leží čtvrtina dat (medián dolní poloviny), \textbf{horní kvartil} $Q_3$ hodnota, pod kterou leží tři čtvrtiny dat (medián horní poloviny). Medián je druhý kvartil. \textbf{Mezikvartilové rozpětí} je $Q_3 - Q_1$.

  \textbf{Krabicový graf} znázorňuje minimum, $Q_1$, medián, $Q_3$ a maximum. Krabice obsahuje prostřední polovinu dat.
\end{definitionbox}

\begin{center}
  \begin{tikzpicture}[scale=1.5]
    \draw[->] (0.5,0) -- (5.5,0);
    \foreach \x in {1,2,3,4,5} \draw (\x,0.06) -- (\x,-0.06) node[below] {\small $\x$};
    \draw[thick] (1,0.6) -- (2,0.6);
    \draw[thick] (3,0.6) -- (5,0.6);
    \draw[thick, fill=colorofwisdom!30] (2,0.35) rectangle (3,0.85);
    \draw[very thick, colorofwisdom] (2,0.35) -- (2,0.85);
    \draw[thick] (1,0.45) -- (1,0.75);
    \draw[thick] (5,0.45) -- (5,0.75);
    \node[above] at (1,0.8) {\tiny min};
    \node[above] at (2,0.9) {\tiny $Q_1 = \tilde{x}$};
    \node[above] at (3,0.9) {\tiny $Q_3$};
    \node[above] at (5,0.8) {\tiny max};
  \end{tikzpicture}
\end{center}

Pro naše známky je dolní polovina $1, 1, 2, 2, 2$ s mediánem $Q_1 = 2$ a horní polovina $2, 3, 3, 4, 5$ s mediánem $Q_3 = 3$. Medián je také $2$, takže splývá s dolním kvartilem. To je u malých souborů s mnoha stejnými hodnotami běžné.

\begin{problem}
  Výšky (v cm) osmi hráčů basketbalového týmu jsou $185, 192, 178, 201, 188, 195, 183, 190$.
  \begin{enumerate}[label=(\alph*)]
    \item Určete průměr, medián a variační rozpětí.
    \item Určete směrodatnou odchylku.
    \item Do týmu přibude hráč s výškou $230$ cm. Jak se změní průměr a medián?
  \end{enumerate}
\end{problem}

\begin{solution}
  \begin{enumerate}[label=(\alph*)]
    \item Součet $1512$, průměr $\bar{x} = 189$. Seřazeno: $178, 183, 185, 188, 190, 192, 195, 201$, medián $\frac{188 + 190}{2} = 189$. Rozpětí $201 - 178 = 23$.
    \item Odchylky: $-11, -6, -4, -1, 1, 3, 6, 12$. Čtverce: $121, 36, 16, 1, 1, 9, 36, 144$, součet $364$. $s^2 = 45{,}5$, $s\approx 6{,}75\ \si{cm}$.
    \item Průměr $\frac{1742}{9}\approx 193{,}6$, vzroste skoro o pět centimetrů. Medián je teď pátá hodnota, tedy $190$, vzroste jen o centimetr. Opět vidíme, jak je průměr citlivý na extrémní hodnoty.
  \end{enumerate}
\end{solution}

\section{Závislost dvou znaků}\label{sec:zavislost-znaku}

Často nás zajímá, jestli spolu dva znaky souvisí. Mají vyšší lidé větší boty? Mají studenti, kteří víc spí, lepší známky? Dvojice hodnot $(x_i, y_i)$ můžeme zakreslit jako body do roviny (\textbf{bodový graf}) a podívat se, zda tvoří nějaký trend.

\begin{definitionbox}
  \textbf{Korelační koeficient} $r\in[-1, 1]$ měří míru \textbf{lineární} závislosti dvou znaků:
  \begin{equation*}
    r = \frac{\sum{(x_i - \bar{x})(y_i - \bar{y})}}{\sqrt{\sum{(x_i - \bar{x})}^2}\cdot\sqrt{\sum{(y_i - \bar{y})}^2}}.
  \end{equation*}
  \begin{itemize}
    \item $r$ blízko $1$: silná přímá závislost (s rostoucím $x$ roste $y$),
    \item $r$ blízko $-1$: silná nepřímá závislost,
    \item $r$ blízko $0$: žádná lineární závislost (ale může existovat jiná, třeba kvadratická).
  \end{itemize}
\end{definitionbox}

\begin{center}
  \begin{tikzpicture}[scale=0.8]
    \foreach \sx/\pts/\lab in {
      0/{0.2/0.28, 0.44/0.68, 0.67/0.74, 0.91/1.1, 1.15/1.31, 1.38/1.1, 1.62/1.26, 1.85/2.02, 2.09/1.8, 2.33/1.98, 2.56/2.69, 2.8/2.52}/{$r\approx 0{,}96$},
      3.7/{0.2/2.62, 0.44/2.26, 0.67/2.17, 0.91/1.74, 1.15/1.8, 1.38/1.75, 1.62/1.4, 1.85/1.33, 2.09/1.12, 2.33/0.63, 2.56/0.81, 2.8/0.55}/{$r\approx -0{,}98$},
      7.4/{0.2/2.05, 0.44/2.11, 0.67/1.37, 0.91/0.88, 1.15/0.3, 1.38/1.76, 1.62/1.33, 1.85/1.97, 2.09/1.12, 2.33/1.99, 2.56/0.9, 2.8/2.06}/{$r\approx 0$},
      11.1/{0.2/0.73, 0.44/1.17, 0.67/1.65, 0.91/2.05, 1.15/2.48, 1.38/2.47, 1.62/2.51, 1.85/2.32, 2.09/2.12, 2.33/1.71, 2.56/1.23, 2.8/0.66}/{$r\approx 0$}} {
      \begin{scope}[xshift=\sx cm]
        \draw[->] (0,0) -- (3.1,0) node[right] {\scriptsize $x$};
        \draw[->] (0,0) -- (0,3) node[above] {\scriptsize $y$};
        \foreach \x/\y in \pts \filldraw[colorofwisdom] (\x,\y) circle (2.2pt);
        \node[font=\small] at (1.5,-0.45) {\lab};
      \end{scope}
    }
    \node[font=\scriptsize] at (1.5,-0.95) {přímá závislost};
    \node[font=\scriptsize] at (5.2,-0.95) {nepřímá závislost};
    \node[font=\scriptsize] at (8.9,-0.95) {žádná závislost};
    \node[font=\scriptsize, align=center] at (12.6,-1.1) {silná, ale ne lineární!};
  \end{tikzpicture}
\end{center}
Korelační koeficient si můžete představit takhle. Bodovým grafem proložíme svislou a vodorovnou přímku průměry $\bar{x}$ a $\bar{y}$, čímž ho rozdělíme na čtyři kvadranty. Body vpravo nahoře a vlevo dole přispívají do čitatele kladně (obě odchylky mají stejné znaménko), body v ostatních dvou kvadrantech záporně. Pokud body tvoří stoupající pruh, převažují kladné příspěvky a $r > 0$. Na posledním obrázku se příspěvky z levé a pravé poloviny paraboly přesně vyruší, takže $r\approx 0$, přestože $y$ je na $x$ silně závislé. Proto je dobré se vždy nejdřív podívat na graf.

Pozorný čtenář si všimne, že vzorec je přesně kosinus úhlu mezi vektory odchylek $(x_1 - \bar{x}, \ldots, x_n - \bar{x})$ a $(y_1 - \bar{y}, \ldots, y_n - \bar{y})$ podle kapitoly~\ref{kap:analyticka}, jen v $n$-rozměrném prostoru. Proto je vždy mezi $-1$ a $1$.

\begin{note}
  \textbf{Korelace neznamená kauzalitu.} V létě roste prodej zmrzliny i počet utonulých. Zmrzlina ale k utonutí nevede. Obojí způsobuje teplé počasí. Na internetu najdete stránky se sbírkami absurdních korelací. Například počet filmů, ve kterých hrál Nicolas Cage, v letech 1999 až 2009 překvapivě dobře koreloval s počtem lidí, kteří utonuli v bazénu. Pokud někdo tvrdí, že `$A$ způsobuje $B$', protože spolu korelují, vždy se zeptejte, jestli za obojím nestojí nějaké $C$.
\end{note}

\section{Příběh --- Účty smrti a růže z Krymu}

Moderní statistika se narodila mezi mrtvými. V Londýně se od 16. století každý týden vydávaly takzvané \textit{Bills of Mortality}, tedy seznamy úmrtí podle farností a příčin. Pomáhaly sledovat, kdy se blíží další vlna moru. Mezi příčinami smrti v nich najdete například `zemřel leknutím', `planetou zasažen' nebo `šílenství'. Obchodník John Graunt je v roce 1662 poprvé systematicky zpracoval. Zjistil, že se rodí trochu víc chlapců než dívek, že úmrtnost dětí je děsivě vysoká a že čísla se dají použít k odhadu počtu obyvatel. Jeho kniha je považována za počátek demografie i statistiky vůbec.

O dvě století později se do Krymské války vydala ošetřovatelka Florence Nightingaleová. Ve vojenských nemocnicích zjistila, že vojáci umírají mnohem častěji na nemoci jako cholera nebo tyfus než na zranění z bitev. Aby o tom přesvědčila politiky, kteří čísla nečetli, vymyslela nový druh grafu, takzvaný \textit{coxcomb} (polární diagram, někdy nazývaný Nightingaleova růže). Každý měsíc tvořil jednu výseč a plocha barevných částí ukazovala, kolik vojáků zemřelo na nemoci, na zranění a z jiných příčin. Graf byl tak výmluvný, že vedl k reformě vojenských nemocnic a hygieny. Florence Nightingaleová se v roce 1858 stala první ženou zvolenou do britské Královské statistické společnosti. Lidé si ji pamatují jako `dámu s lampou'. Měli by si ji pamatovat i jako dámu s grafem.

\section{Shrnutí a otázky k ústní zkoušce}

\begin{problem}
  Odpovězte vlastními slovy:
  \begin{enumerate}
    \item Co je statistický soubor, jednotka a znak? Jaké druhy znaků znáte?
    \item Co je absolutní, relativní a kumulativní četnost?
    \item Jaké znáte způsoby grafického znázornění dat? Jak lze grafem manipulovat?
    \item Definujte aritmetický průměr, medián a modus. Kdy je lepší použít medián?
    \item Kdy se používá geometrický a harmonický průměr? Uveďte příklad.
    \item Co je rozptyl a směrodatná odchylka? Proč se v rozptylu odchylky umocňují?
    \item Co jsou kvartily a krabicový graf?
    \item Co měří korelační koeficient? Proč korelace neznamená kauzalitu?
  \end{enumerate}
\end{problem}
